\section{Límites del sesgo estructural: ViT, Cross-attention, Fourier e ImageNet}

Con el latent bottleneck establecido, esta sección responde dos preguntas: ¿en qué se distingue Perceiver de ViT y de la cross-attention ya existente? ¿Y cómo obtiene la posición sin depender de convoluciones bidimensionales? La exposición no sostiene que toda información estructural previa sea inútil; mediante comparaciones y ablaciones busca la información de entrada «mínima pero suficiente».

\subsection{Comparación con ViT: el patching ya es una suposición sobre imágenes}

La sección anterior presentó el cuello de botella unificado de Perceiver; esta sección lo compara primero con el Transformer visual más cercano. ViT divide la imagen en patches bidimensionales fijos y después genera tokens mediante una convolución o una proyección lineal. Este enfoque es muy eficaz para imágenes, pero exige redefinir el patching para cada tipo de cuadrícula, nube de puntos o sensor. Perceiver puede recibir directamente píxeles o elementos de un arreglo y trata la posición como una característica adicional.

\lecturefigure{slide-12-contrast-vit.jpg}{ViT reduce el número de tokens mediante un patch embedding bidimensional; Perceiver evita incorporar los patches de cuadrícula a su interfaz central.}{Video oficial actual de Stanford Online, 00:21:42--00:22:55.}

\begin{knowledgebox}{Lectura de la figura: ViT y Perceiver optimizan cuellos de botella distintos}
ViT reduce primero $M$ mediante el tamaño de patch y después sigue aplicando self-attend sobre los patch tokens; Perceiver conserva una entrada más fina y controla el workspace mediante latent queries. El primero suele ser más eficiente y está más consolidado; el segundo se traslada con mayor facilidad a arreglos sin cuadrícula o heterogéneos.
\end{knowledgebox}

\teachervoice{En la sesión de preguntas y respuestas, el docente fija una postura clara: en visión, los convolution/attention hybrids buscan un punto de Pareto preciso entre velocidad y calidad para problemas de imágenes, y eso es muy importante; el equipo de Perceiver busca un modelo tan general como sea posible que conserve un desempeño utilizable. Las dos líneas no se excluyen mutuamente.}

\subsection{La cross-attention no es exclusiva de Perceiver}

DETR usa object queries para decodificar bounding boxes a partir de mapas de características convolucionales; Slot Attention usa learned slots que leen las características de forma competitiva y forman representaciones centradas en objetos. La novedad de Perceiver consiste en elevar este tipo de attention asimétrica a un cuello de botella de entrada de uso general y procesarla repetidamente en el latent workspace.

\lecturefigure{slide-13-cross-attention-detr.jpg}{Precedentes de cross-attention en trabajos de visión como DETR y Slot Attention.}{Video oficial actual de Stanford Online, 00:22:56--00:23:11.}

\begin{importantbox}{La atribución del método debe ser precisa}
Perceiver no inventó la cross-attention. Su aportación es una combinación: la lectura asimétrica de una entrada grande hacia un learned latent pequeño, el procesamiento profundo dentro del latent, las suposiciones débiles sobre la modalidad y, en Perceiver IO, la decodificación mediante consultas de salida.
\end{importantbox}

\subsection{Fourier positional features: representar coordenadas con funciones base de frecuencia}

Para expresar la posición con detalle fino sin codificar de forma fija la conectividad local, esta sección construye, para la coordenada bidimensional $p=(x,y)$, características seno/coseno de varias frecuencias:

\[
\gamma(p)=\big[\sin(2\pi f_1x),\cos(2\pi f_1x),\sin(2\pi f_1y),\cos(2\pi f_1y),\ldots\big].
\]

Aquí, $f_k$ es un conjunto de frecuencias que va de las bajas hasta las cercanas a la de Nyquist; $x,y$ son coordenadas normalizadas; $\gamma(p)$ se concatena con el RGB u otras características de contenido. Las frecuencias bajas expresan la posición a gran escala y las altas ayudan a distinguir coordenadas finas.

\lecturefigure{slide-14-fourier-position-encodings.jpg}{Fourier positional encodings bidimensionales: características de coordenadas de varias frecuencias concatenadas con el RGB.}{Explicación original, 00:26:18--00:26:55; la progressive build completa vuelve a mostrarse en el video oficial actual, 00:39:52.}

\teachervoice{El docente reporta una observación concreta: con entradas de imagen, hacer \textbf{concatenate} de la Fourier position con el RGB suele funcionar mejor que proyectar primero y luego sumar. No da un mecanismo definitivo; solo conjetura que las características de contenido de una imagen no son tan dispersas como un embedding de palabras. Conviene conservar este estado de «observación confiable, explicación pendiente».}

\begin{warningbox}{Las Fourier features siguen aportando estructura real}
Aunque el modelo no tiene vecindades convolucionales, la codificación de coordenadas le indica explícitamente en qué lugar del plano bidimensional está cada píxel. Decir que «no usa en absoluto la estructura de la imagen» es inexacto; es más preciso decir que la estructura se proporciona como característica de entrada y no como un grafo de cómputo fijo.
\end{warningbox}

\subsection{Permuted ImageNet: la invariancia a permutaciones no implica que la estructura sea inútil}

A continuación, un experimento contraintuitivo comprueba si las características de posición realmente sustituyen al orden fijo del arreglo. El experimento reordena al azar los píxeles de la imagen; mientras cada píxel conserve su codificación de posición correcta, Perceiver obtiene resultados top-1 en ImageNet similares. Este resultado muestra que la arquitectura no depende del orden en memoria del arreglo de entrada ni de un orden de recorrido local; no demuestra que la posición, la localidad o el sesgo convolucional carezcan de valor para la eficiencia en datos.

\lecturefigure{slide-15-imagenet-permuted-classification.jpg}{ImageNet con entrada estándar y permuted: las características de posición correctas hacen que el modelo no dependa de la permutación del arreglo.}{Video oficial actual de Stanford Online, 00:29:34--00:33:47.}

\begin{knowledgebox}{Lectura de la tabla: se compara el orden del arreglo, no la eliminación del espacio}
La versión permuted sigue incluyendo la codificación de posición, por lo que el modelo puede aprender qué píxeles son vecinos en el espacio. Si también se revolvieran las etiquetas de posición, la tarea se convertiría en un problema distinto. Los resultados de la tabla respaldan la input-order robustness, no la idea de que «la estructura bidimensional no existe».
\end{knowledgebox}

\teachervoice{La advertencia importante de la sesión de preguntas y respuestas es que un modelo con sesgo débil puede necesitar más datos para recuperar la estructura por sí mismo. Con conjuntos grandes como ImageNet, la generalidad puede obtenerse mediante entrenamiento; el escenario con pocos datos sigue siendo un cuello de botella sin resolver, como reconoce explícitamente el expositor.}

\subsection{Resumen de la sección}

Perceiver incorpora menos suposiciones de patch/cuadrícula que ViT, pero sigue usando características de posición. La cross-attention tiene numerosos precedentes; la aportación está en la combinación de uso general. Permuted ImageNet muestra robustez frente al orden de la entrada, pero no anula la ventaja en eficiencia de datos del sesgo local.
